\section{Representaciones y aplicaciones adicionales}

Hasta ahora, muchos de los estudios sobre autómatas celulares neuronales se han enfocado en la morfogénesis, que es el proceso de formación de formas, y en su habilidad para crear estructuras visuales que se organizan solas. Sin embargo, también han surgido enfoques que amplían el marco conceptual hacia nuevas formas de representación y aplicaciones en dominios específicos. A continuación, se presentan algunas de estas variantes.

Mattie Tesfaldet et al. \cite{ViTCA} extienden los autómatas celulares neuronales mediante mecanismos de atención inspirados en arquitecturas de \textit{transformers}, dando lugar a Autómatas Celulares con Transformadores de Visión. Esta arquitectura híbrida permite que las interacciones locales tengan acceso contextual a regiones más amplias a través de \textit{self-attention}, lo que habilita organización global sin incurrir en la complejidad cuadrática típica de la atención completa. Autómatas Celulares con Transformadores de Visión adapta el proceso de los transformadores de visión tradicionales, que dividen una imagen en parches y los convierten en secuencias de \textit{tokens}, al formato de una cuadrícula de celdas, donde cada celda almacena información visual, posicional y oculta. En lugar de aplicar atención global, restringe el cálculo a vecindarios locales, manteniendo así la naturaleza distribuida del sistema y reduciendo el costo computacional. El entrenamiento se realiza de forma iterativa, alternando entre inyectar entradas corruptas y actualizar parcialmente los estados celulares hasta reconstruir la imagen. De este modo, Autómatas Celulares con Transformadores de Visión combina eficiencia, capacidad de autoorganización y organización espacial coherente.

En un ámbito distinto, Jian Tang et al. \cite{JianNCA} emplean red neuronal convolucional para simular el desarrollo de microestructuras en la solidificación de metales. Cada celda del dominio contiene un estado de fase $P\in[0,1]$ (donde $p=0$ representa líquido y $p=1$ sólido), tres ángulos de Euler $\theta=(\alpha_1,\alpha_2,\alpha_3)$ que definen su orientación cristalográfica, y la magnitud de subenfriamiento $\Delta T=T_{melting}-T$. La evolución se modela con un autómata celular neuronal que aplica recurrentemente una red neuronal convolucional para predecir los incrementos en fase, orientación y estados ocultos. Estos canales ocultos facilitan el aprendizaje de patrones temporales complejos sin necesidad de reglas explícitas. Con ello, el modelo reproduce fenómenos como crecimiento competitivo o direcciones preferentes, con una eficiencia computacional sustancialmente mayor que la de muchos autómatas celulares tradicionales y con capacidad de generalización fuera del rango de entrenamiento.

Por otro lado, Gordon H. Y. Li et al. \cite{PNCA} proponen el modelo autómata celular neuronal fotónico, que aprovecha la velocidad y la interconexión masiva de la fotónica para implementar autómata celular neuronal en hardware óptico. A diferencia de arquitecturas convencionales basadas en perceptrón multicapa o red neuronal convolucional, difíciles de trasladar a entornos fotónicos por sus conexiones densas, autómata celular neuronal fotónico utiliza reglas locales aprendidas mediante técnicas de aprendizaje profundo. Esto permite realizar clasificación de imágenes en sistemas ópticos, donde la actualización celular se implementa mediante interferencia lineal de luz y activaciones ópticas no lineales. Los autores demuestran clasificación binaria con alta precisión utilizando un número muy reducido de parámetros fotónicos programables, mostrando además robustez frente a datos fuera de distribución. El autómata celular neuronal fotónico representa una vía prometedora para llevar el aprendizaje autoorganizado a plataformas fotónicas escalables.

En el ámbito de la biología computacional, Caitlin Grasso y Josh Bongard \cite{ENCA} introducen los Autómatas Celulares Neuronales Potenciados. Esta es una extensión que incorpora el concepto de \textit{empowerment} —la cantidad de control que una célula puede ejercer sobre su entorno— formalizado mediante medidas de información mutua entre acciones pasadas y estados sensoriales futuros. Incluir \textit{empowerment} como un objetivo adicional en las tareas de morfogénesis mejora la capacidad de evolución y genera patrones de crecimiento más fuertes y adaptables. Esto permite que cada célula actúe como un agente independiente dentro de un sistema distribuido y facilita el estudio de la coordinación y comunicación celular desde una perspectiva computacional.

Finalmente, mientras que modelos como Autómatas Celulares con Transformadores de Visión, Autómata Celular Neuronal Fotónico o Autómatas Celulares Neuronales Potenciados amplían las capacidades arquitectónicas y las plataformas de aplicación de los autómatas celulares neuronales, Richardson et al. \cite{richardsonNCA} utilizan los autómatas celulares neuronales como herramienta para modelar fenómenos dinámicos complejos. Entrenando autómata celular neuronal mediante optimización basada en gradientes para reproducir trayectorias espacio-temporales generadas por ecuaciones diferenciales parciales (por ejemplo, Gray–Scott), este enfoque captura tanto estructuras transitorias como patrones estacionarios —incluyendo la formación de patrones tipo Turing— y aborda el problema inverso de inferir reglas locales a partir de comportamiento emergente. Con una complejidad de red mínima, los autores muestran capacidad de generalización, respeto por simetrías y robustez frente a perturbaciones, posicionando a los autómatas celulares neuronales como un marco explicativo y eficiente para sistemas biológicos y físicos.

